\documentclass{article}
\usepackage{lmodern}
\usepackage{amsmath}
\usepackage{listings}
\usepackage{fullpage}
\usepackage{parskip}
\usepackage{xcolor}
\lstset{
	basicstyle=\ttfamily,
	columns=fullflexible,
	frame=single,
	breaklines=true,
	postbreak=\mbox{\textcolor{red}{$\hookrightarrow$}\space},
}
\begin{document}
\title{Experiment `p\_4\_1\_d\_sum\_odd\_numbers' Results}
\author{\today}
\date{}
\maketitle
\textbf{Experiment outcome: }SUCCESS
\\\textbf{Bad responses: }0
\\\textbf{Responses containing}~\texttt{assume}~\textbf{: }0
\\\textbf{Responses containing}~\texttt{expect}~\textbf{: }0
\\\textbf{Resolution attempts: }1
\\\textbf{Hard fails (resolution): }0
\\\textbf{Soft fails (resolution): }0
\\\textbf{Verification attempts: }1
\section*{Problem Specification}
\textbf{Problem name: }p\_4\_1\_d\_sum\_odd\_numbers
\\\textbf{Natural language statement: }null
\\\textbf{Method signature: }p\_4\_1\_d\_sum\_odd\_numbers(a: int, b: int) returns (sum: int)
\subsection*{Ensures}
\begin{itemize}
\item \texttt{sum == sum\_odd(a, b)}
\end{itemize}
\subsection*{Functional Code Given}
\begin{lstlisting}
function sum_odd(a: int, b: int): int
decreases b - a
{
  if a > b then 0
  else if a % 2 == 1 then a + sum_odd(a + 1, b)
  else sum_odd(a + 1, b)
}
\end{lstlisting}
\clearpage
\section*{GenAI interactions}
Below you will find all interactions between the `user' (program) and the `assistant' (OpenAI).
\subsection*{Program $\rightarrow$ GenAI}
\begin{lstlisting}
You are given the following task to perform in Dafny:

Write a method to satisfy the below requirements.

The signature should be:

p_4_1_d_sum_odd_numbers(a: int, b: int) returns (sum: int)

The method should respect the following contract:

ensures sum == sum_odd(a, b)

The contract uses the following dafny code:

function sum_odd(a: int, b: int): int
decreases b - a
{
  if a > b then 0
  else if a % 2 == 1 then a + sum_odd(a + 1, b)
  else sum_odd(a + 1, b)
}

These function(s) must not be used in your implementation of the method.

Produce and show only the Dafny body of this method, including the curly braces that surround it. Do not show the signature nor contract. You must not use 'assume' nor 'expect' anywhere in your code.
\end{lstlisting}
\subsection*{GenAI $\rightarrow$ Program}
\textbf{System fingerprint: }null
\\\textbf{ID: }chatcmpl-DnLFYhoDcVepiRNr1Wg23Gx1X62av
\\\textbf{Model: }gpt-5.5-2026-04-23
\\\textbf{Created at: }1780652128
\\\textbf{Finish reason: }stop
\\\textbf{Usage: }444 tokens in, and 3881 tokens out
\begin{lstlisting}
{
  if a > b {
    sum := 0;
    assert sum_odd(a, b) == 0;
  } else {
    sum := 0;
    var i := a;
    while i <= b
      invariant a <= i <= b + 1
      invariant sum + sum_odd(i, b) == sum_odd(a, b)
      decreases b - i + 1
    {
      if i % 2 == 1 {
        assert sum_odd(i, b) == i + sum_odd(i + 1, b);
        sum := sum + i;
      } else {
        assert sum_odd(i, b) == sum_odd(i + 1, b);
      }
      assert sum + sum_odd(i + 1, b) == sum_odd(a, b);
      i := i + 1;
    }
    assert i == b + 1;
    assert sum_odd(i, b) == 0;
    assert sum == sum_odd(a, b);
  }
}
\end{lstlisting}
\section*{Final Program}
\begin{lstlisting}
function sum_odd(a: int, b: int): int
decreases b - a
{
  if a > b then 0
  else if a % 2 == 1 then a + sum_odd(a + 1, b)
  else sum_odd(a + 1, b)
}

method p_4_1_d_sum_odd_numbers(a: int, b: int) returns (sum: int)
	ensures sum == sum_odd(a, b)
{
  if a > b {
    sum := 0;
    assert sum_odd(a, b) == 0;
  } else {
    sum := 0;
    var i := a;
    while i <= b
      invariant a <= i <= b + 1
      invariant sum + sum_odd(i, b) == sum_odd(a, b)
      decreases b - i + 1
    {
      if i % 2 == 1 {
        assert sum_odd(i, b) == i + sum_odd(i + 1, b);
        sum := sum + i;
      } else {
        assert sum_odd(i, b) == sum_odd(i + 1, b);
      }
      assert sum + sum_odd(i + 1, b) == sum_odd(a, b);
      i := i + 1;
    }
    assert i == b + 1;
    assert sum_odd(i, b) == 0;
    assert sum == sum_odd(a, b);
  }
}
\end{lstlisting}
\section*{Total Token Usage}
\textbf{Input tokens: }444
\\\textbf{Output tokens: }3881
\\\textbf{Reasoning tokens: }3581
\\\textbf{Sum of `total tokens': }4325
\section*{Experiment Timings}
\textbf{Overall Experiment} started at 1780652127990, ended at 1780652233621, lasting 105631ms (105.63 seconds)
\\\textbf{Iteration \#1} started at 1780652127996, ended at 1780652233621, lasting 105625ms (105.63 seconds)
\end{document}
